\providecommand{\Res}[2]{\qty[round-mode=figures,round-precision=3]{\fpeval{#1}}{#2}}
\providecommand{\Val}[1]{\num[round-mode=figures,round-precision=4]{\fpeval{#1}}}
\providecommand{\ResSci}[2]{\qty[round-mode=figures,round-precision=3,exponent-mode=scientific]{\fpeval{#1}}{#2}}
\subsubsection*{Question a) Résistances}
\begin{enumerate}[label=\arabic*)]
    \item \textbf{Variables connues et recherchées :} $U = \qty{120}{V}$, $P_1 = \qty{\v{P1}}{W}$, $P_2 = \qty{\v{P2}}{W}$ ; variables recherchées : $R_1$, $R_2$
    \item \textbf{Équations pertinentes :}
    \[ P = UI \qquad U = RI \quad\Rightarrow\quad P = \frac{U^2}{R} \]
    \item \textbf{Isolation de la variable recherchée :}
    \[ R = \frac{U^2}{P} \]
    \item \textbf{Résolution numérique :}
    \[ R_1 = \frac{(\qty{120}{V})^2}{\qty{\v{P1}}{W}} = \Res{120^2/\v{P1}}{\ohm} \qquad R_2 = \frac{(\qty{120}{V})^2}{\qty{\v{P2}}{W}} = \Res{120^2/\v{P2}}{\ohm} \]
    \item \textbf{Réponse :} $R_1 = \mathbf{\Res{120^2/\v{P1}}{\ohm}}$ et $R_2 = \mathbf{\Res{120^2/\v{P2}}{\ohm}}$.
\end{enumerate}

\subsubsection*{Question b) Courant en série}
\[ I = \frac{U}{R_1 + R_2} = \frac{\qty{120}{V}}{\Val{120^2/\v{P1}+120^2/\v{P2}}\,\unit{\ohm}} = \Res{120/(120^2/\v{P1}+120^2/\v{P2})}{A} \]
Le courant dans le circuit est de \textbf{\Res{120/(120^2/\v{P1}+120^2/\v{P2})}{A}}.

\subsubsection*{Question c) Puissances réelles}
\[ P_1' = R_1 I^2 = \Res{120^2/\v{P1}*(120/(120^2/\v{P1}+120^2/\v{P2}))^2}{W} \qquad P_2' = R_2 I^2 = \Res{120^2/\v{P2}*(120/(120^2/\v{P1}+120^2/\v{P2}))^2}{W} \]
En série, le courant est le même : l'ampoule de plus grande résistance, celle de \qty{\v{P1}}{W}, dissipe le plus de puissance et \textbf{brille le plus}, contrairement à ce qui se passe quand chacune est branchée seule sur \qty{120}{V}.
